\chapter{Interface Utilisateur et Mise en Œuvre}
\label{ch:ihm}

\section{Téléopération par manette Xbox}
\label{sec:ihm_teleop}

La téléopération repose sur une manette sans fil (compatible Xbox), pilotée par
\texttt{joy\_node}/\texttt{teleop\_twist\_joy} (\texttt{joystick\_drive.launch.py}), et non
le \texttt{teleop\_twist\_keyboard} suggéré au cahier des charges.

\begin{table}[H]
	\centering
	\caption{Configuration de la téléopération manette (\texttt{xbox\_teleop.yaml}).}
	\label{tab:teleop_params}
	\begin{tabularx}{\linewidth}{L{5.5cm} C{2.6cm} L{4.3cm}}
		\toprule
		\textbf{Paramètre} & \textbf{Valeur} & \textbf{Remarque} \\
		\midrule
		\texttt{require\_enable\_button} & \texttt{false} & aucune touche « homme-mort » requise \\
		Axe linéaire / angulaire & axe 1 (éch. 0.25) / axe 3 (éch. 2.4) & sticks $\to$ \topic{/cmd\_vel} \\
		Zone morte / répétition & 0.08 / 20\,Hz & \texttt{joy\_node} \\
		\bottomrule
	\end{tabularx}
\end{table}

L'absence de bouton d'activation obligatoire signifie que toute déviation accidentelle des
manches se traduit immédiatement par un mouvement, sans touche à maintenir --- une marge de
sécurité réduite par rapport à un dispositif « homme-mort » classique
(chapitre~\ref{ch:limitations}). \texttt{joystick\_drive.launch.py} démarre conjointement
\texttt{joy\_node}, \texttt{teleop\_twist\_joy\_node} et \texttt{motor\_bridge\_node}
(vitesse angulaire max. relevée à 3.0\,rad/s, rapport cyclique minimal 22\,\%).

\section{Fusion \gls{ekf} --- détail de configuration}
\label{sec:ihm_ekf_detail}

\texttt{ekf\_odom\_imu.yaml} (\texttt{ekf\_fusion.launch.py}) configure
\texttt{ekf\_filter\_node} à \textbf{30\,Hz}, \texttt{two\_d\_mode}~: de l'odométrie
(\topic{/odom}), seules les \emph{vitesses} $(v_x,v_y,\dot\theta)$ sont fusionnées (la
position absolue n'est pas injectée, évitant un double comptage de dérive) ; de l'\gls{imu}
(\topic{/imu/data\_raw}), l'orientation de lacet absolue, $\dot\theta$ et l'accélération
longitudinale $a_x$ sont retenues. Le \texttt{world\_frame} reste \texttt{odom}, laissant
la correction \topic{map}$\to$\topic{odom} à \texttt{slam\_toolbox}/\gls{amcl} --- une
séparation propre entre estimation locale (\gls{ekf}) et globale (\gls{slam}/\gls{amcl}).

\section{Tableau de bord web de télémétrie}
\label{sec:ihm_dashboard}

Le nœud \texttt{imu\_web\_dashboard} (\texttt{esibot\_dashboard}) expose un serveur HTTP
autonome (\texttt{http.server}, sans dépendance externe), diffusant une page HTML/JS à page
unique, sensiblement différente du tableau de bord du cahier des charges (qui prévoyait une
passerelle WebSocket \texttt{rosbridge\_suite} avec carte, vidéo et boutons de
téléopération) --- l'implémentation réelle se concentre sur la \textbf{télémétrie}.

\begin{table}[H]
	\centering
	\caption{Widgets du tableau de bord \texttt{imu\_web\_dashboard}.}
	\label{tab:dashboard_widgets}
	\begin{tabularx}{\linewidth}{L{3.2cm} L{8.7cm}}
		\toprule
		\textbf{Widget} & \textbf{Contenu} \\
		\midrule
		Orientation & roulis/tangage (intégration accéléromètre), horizon artificiel \texttt{<canvas>}, mise à zéro (\texttt{POST /api/zero}) \\
		Odométrie & position/cap/vitesses (\topic{/odometry/filtered}), trajectoire (1200 pts) \\
		Manette \& commande & axes/boutons \topic{/joy} et conversion \topic{/cmd\_vel} $\to$ vitesses de roue (équation~\ref{eq:diff_drive}) \\
		\bottomrule
	\end{tabularx}
\end{table}

L'interface interroge en boucle \texttt{GET /api/imu} (état JSON consolidé, protégé par
verrou pour la cohérence entre thread \ros{}\,2 et thread HTTP). Le tableau de bord
\textbf{n'affiche ni la carte \gls{slam} ni le flux vidéo}, deux fonctionnalités du
Livrable~5 (chapitre~\ref{ch:limitations}).

\section{Flux vidéo ESP32-CAM}
\label{sec:ihm_camera}

Le pont \texttt{esp32\_camera\_bridge} (\texttt{mjpeg\_image\_publisher}) se connecte en
client HTTP au flux \gls{mjpeg} de l'ESP32-CAM, extrait chaque trame JPEG
(marqueurs \texttt{0xFFD8}/\texttt{0xFFD9}), la décode avec OpenCV et republie sur
\topic{/esp32/camera/image\_raw} via \texttt{cv\_bridge}, avec reconnexion automatique
(2\,s). Ce flux est consommable mais non consommé par un nœud de vision dédié --- aucun
\texttt{vision\_node} (détection ligne/obstacles, prévu au cahier des charges) n'a été
retrouvé dans le code livré.

\section{Mise en Œuvre et Procédures d'Exploitation}
\label{ch:exploitation}

Le dépôt fournit trois scripts shell (\texttt{scripts/}) qui illustrent la séparation
bas-niveau/haut-niveau, contrairement au test T1 du cahier des charges qui suppose un
unique \texttt{esibot\_full.launch.py}.

\begin{table}[H]
	\centering
	\caption{Scripts de démarrage disponibles dans \texttt{scripts/}.}
	\label{tab:startup_scripts}
	\begin{tabularx}{\linewidth}{L{3.6cm} L{4.6cm} L{4.6cm}}
		\toprule
		\textbf{Script} & \textbf{Rôle} & \textbf{Particularité} \\
		\midrule
		\texttt{start\_esibot.sh} & lancement manuel du calibrateur & script de développement minimal \\
		\texttt{start\_esibot\_all.sh} & démarrage production (téléop, \gls{imu}, \gls{ekf}, dashboard) & exige une calibration valide ; radar désactivé \\
		\texttt{start\_esibot\_all\_no\_\allowbreak calib.sh} & variante sans vérification & active le radar ultrasonique au lieu du Kinect \\
		\bottomrule
	\end{tabularx}
\end{table}

\begin{notebox}[title={Aucun script ne démarre le pipeline Kinect/\gls{slam}}]
	Les trois scripts se limitent à la couche « commande bas-niveau + télémétrie »
	(chapitre~\ref{ch:commande_bas_niveau}). Le pipeline de cartographie
	(\texttt{mapping.launch.py}) et Nav2 (\texttt{navigation2.launch.py}) doivent être
	démarrés \emph{séparément}, typiquement depuis le PC opérateur (couche haut-niveau,
	figure~\ref{fig:global_archi}) --- reflet direct du découpage en trois couches
	(chapitre~\ref{ch:archi_globale}).
\end{notebox}

Avant tout démarrage via \texttt{start\_esibot\_all.sh}, le script vérifie la validité
d'un fichier de calibration JSON (\texttt{\textasciitilde/.esibot/pwm\_rpm\_map.json},
généré par \texttt{motor\_calibrator}, section~\ref{sec:cbn_calibrator}), avec coefficients
positifs exigés pour les 4 combinaisons \{gauche,droite\}$\times$\{avant,arrière\} --- un
garde-fou bloquant le démarrage en cas de calibration absente ou invalide.

\begin{lstlisting}[style=BashStyle, caption={Séquence complète --- démarrage bas-niveau, puis cartographie et navigation.}, label={lst:slam_nav_sequence}]
./scripts/start_esibot_all.sh        # teleop + IMU + EKF + dashboard (RPi4)

ros2 launch esibot_slam mapping.launch.py start_rviz:=true        # cartographie
ros2 run nav2_map_server map_saver_cli -f ~/esibot_ws/maps/salle_test
ros2 launch esibot_navigation2 navigation2.launch.py \
    map:=$HOME/esibot_ws/maps/salle_test.yaml use_rviz:=true      # navigation
\end{lstlisting}

Le workspace (\texttt{\textasciitilde/esibot\_ws}) regroupe les huit paquets
\texttt{esibot\_*} et les dépendances tierces \texttt{kinect\_ros2} /
\texttt{esp32\_camera\_bridge}, en complément des paquets standards \ros{}\,2 Humble
(\texttt{slam\_toolbox}, \texttt{nav2\_bringup}, \texttt{robot\_localization},
\texttt{depthimage\_to\_laserscan}, \texttt{joy}, \texttt{teleop\_twist\_joy},
\texttt{cv\_bridge}) --- procédure d'installation détaillée au Livrable~8 (README).
